\documentclass[11pt]{article}
\usepackage[a4paper,margin=1.8cm]{geometry}
\usepackage{ctex}
\usepackage{amsmath,amssymb}
\usepackage{xcolor}
\usepackage{tikz}
\usepackage{tcolorbox}
\usepackage{array}
\usepackage{booktabs}
\usetikzlibrary{arrows.meta,calc}
\tcbuselibrary{skins,breakable}

\definecolor{curve}{RGB}{30,90,180}
\definecolor{secline}{RGB}{220,60,60}
\definecolor{tang}{RGB}{230,140,20}
\definecolor{grayt}{RGB}{110,110,110}
\definecolor{boxbg}{RGB}{245,249,255}

\pagestyle{empty}
\begin{document}

\begin{center}
{\Large\bfseries 非线性函数：某点切线斜率绝对值 $>$ 割线斜率绝对值}\\[2pt]
{\small 几何意义图示\quad(在 $[a,b]$ 连续、在 $(a,b)$ 可导、且非线性)}
\end{center}

\begin{center}
\begin{tikzpicture}[x=2.05cm,y=2.05cm,>=Stealth,font=\small]
  % 坐标轴
  \draw[->,gray] (-0.1,0)--(3.95,0) node[right,black]{$x$};
  \draw[->,gray] (0,-0.1)--(0,1.95) node[above,black]{$f$};

  % 曲线 f(x)=0.4(x-2)^2+0.5
  \draw[curve,very thick,domain=0.6:3.6,samples=100,smooth]
       plot (\x,{0.4*(\x-2)*(\x-2)+0.5});

  % 关键点
  \coordinate (A) at (0.6,1.284);
  \coordinate (B) at (3.6,1.524);
  \coordinate (X) at (3.4,1.284);

  % 割线 AB
  \draw[secline,very thick] (A)--(B);

  % 割线的斜率三角形：run 3, rise 0.24
  \draw[grayt,thin,dashed] (A)--(3.6,1.284);
  \draw[grayt,thin,dashed] (3.6,1.284)--(B);
  \node[grayt,font=\scriptsize,below] at (2.05,1.284){run $=b-a$};

  % 切线 at X, 斜率 1.12
  \draw[tang,very thick] (2.9,0.724)--(3.75,1.676);

  % 切线斜率三角形: run 0.4, rise 0.448
  \draw[tang,thin] (3.0,0.836)--(3.4,0.836)--(3.4,1.284);
  \node[tang,font=\scriptsize,below] at (3.2,0.836){$0.4$};
  \node[tang,font=\scriptsize,right] at (3.4,1.02){$0.448$};

  % 点
  \fill (A) circle (1.8pt) node[below left,black]{$a$};
  \fill (B) circle (1.8pt) node[above left,black]{$b$};
  \fill (X) circle (2pt);
  \draw[dashed,grayt] (X)--(3.4,0) node[below,black]{$\xi$};

  % 数值标签（带边框，避免与曲线重叠混淆）
  \node[secline,draw=secline,fill=white,rounded corners,inner sep=1.5pt,
        font=\scriptsize,anchor=west] at (1.35,1.50){割线斜率 $k=0.08$};
  \node[tang,draw=tang,fill=white,rounded corners,inner sep=1.5pt,
        font=\scriptsize,anchor=east] at (3.35,0.62){切线斜率 $f'(\xi)=1.12$};
\end{tikzpicture}
\end{center}

\vspace{-2mm}
\begin{center}
\small 割线斜率 $k=0.08$；在 $\xi$ 处切线斜率 $f'(\xi)=1.12$，故 $|f'(\xi)|=1.12>|k|=0.08$。
\end{center}

\begin{tcolorbox}[colback=boxbg,colframe=curve!70,breakable,
  title={\bfseries 辅助函数法：先构造 $h$，再分 $k$ 的符号与 $f(c)$ 的位置}]
\small
设 $k=\dfrac{f(b)-f(a)}{b-a}$，\textbf{构造辅助函数}
\[
h(x)=f(x)-f(a)-k(x-a).
\]
则\quad $h(a)=0,\ \ h(b)=f(b)-f(a)-k(b-a)=0$；\quad $h'(x)=f'(x)-k$；\quad
$f$ 非线性 $\iff h\not\equiv 0$。

\vspace{1mm}
\textbf{取偏差非零的内点.} $h$ 连续、$h(a)=h(b)=0$ 且 $h\not\equiv0$，故存在内点 $c\in(a,b)$ 使 $h(c)\neq0$。
记 $L(x)=f(a)+k(x-a)$ 为过 $a,b$ 的弦，则
\[
h(c)=f(c)-L(c)\quad\begin{cases}
>0, & f(c)\ \text{在弦上方},\\[1pt]
<0, & f(c)\ \text{在弦下方}.
\end{cases}
\]

\vspace{1mm}
\textbf{关键恒等式（割线斜率 $=k+$ 偏差/区间长）.}
\[
\frac{f(c)-f(a)}{c-a}=k+\frac{h(c)}{c-a},\qquad
\frac{f(b)-f(c)}{b-c}=k-\frac{h(c)}{b-c}. \tag{$\ast$}
\]
\end{tcolorbox}

\begin{tcolorbox}[colback=white,colframe=tang!80,breakable,
  title={\bfseries 分情况讨论：$k>0$ 或 $k<0$；再分 $f(c)$ 在弦上方或下方}]
\small
在每种情形下，于所取子区间上对 $f(x)-\sigma x$ 用\textbf{罗尔定理}（即拉格朗日中值定理），
得到某 $\xi$ 使 $f'(\xi)=\sigma$，再比较 $|\sigma|$ 与 $|k|$。

\vspace{1mm}
\renewcommand{\arraystretch}{1.35}
\begin{tabular}{>{\centering\arraybackslash}p{1.0cm} >{\centering\arraybackslash}p{2.5cm} >{\centering\arraybackslash}p{1.9cm} >{\centering\arraybackslash}p{4.1cm} >{\centering\arraybackslash}p{3.0cm}}
\toprule
$k$ & $f(c)$ 位置 & 取子区间 & 割线斜率 $\sigma$ & 结论\\
\midrule
$k>0$ & 上方 ($h(c)>0$) & $[a,c]$ & $\sigma=k+\dfrac{h(c)}{c-a}>k>0$ & $|f'(\xi)|=\sigma>|k|$\\
$k>0$ & 下方 ($h(c)<0$) & $[c,b]$ & $\sigma=k-\dfrac{h(c)}{b-c}>k>0$ & $|f'(\xi)|=\sigma>|k|$\\
$k<0$ & 下方 ($h(c)<0$) & $[a,c]$ & $\sigma=k+\dfrac{h(c)}{c-a}<k<0$ & $|f'(\xi)|=|\sigma|>|k|$\\
$k<0$ & 上方 ($h(c)>0$) & $[c,b]$ & $\sigma=k-\dfrac{h(c)}{b-c}<k<0$ & $|f'(\xi)|=|\sigma|>|k|$\\
$k=0$ & $h(c)\neq0$ & $[a,c]\,/\,[c,b]$ & $\sigma=\dfrac{h(c)}{c-a}$ 或 $-\dfrac{h(c)}{b-c}\neq0$ & $|f'(\xi)|>0=|k|$\\
\bottomrule
\end{tabular}

\vspace{0.5mm}
{\footnotesize\textbf{为何这样选区间？} $k>0$ 时：$h(c)>0$ 取\emph{左段}、$h(c)<0$ 取\emph{右段}，都得 $\sigma>k>0$；$k<0$ 时方向相反，取相应一段使 $\sigma<k<0$（即 $|\sigma|>|k|$）。$h(c)\neq0$ 保证严格不等号。\par}
\end{tcolorbox}

\end{document}
